\documentclass[dvipdfmx,a4paper]{article}
\usepackage{amsmath,amssymb,amsthm}
\usepackage[utf8]{inputenc}
\usepackage[dvipdfmx]{geometry}
\geometry{left=25mm,right=25mm,top=25mm,bottom=25mm}

\newtheorem{theorem}{定理}
\newtheorem{definition}{定義}
\newtheorem{lemma}{補題}

\title{二重時空理論における双対ローター遅れの動態方程式\\
――ねじれ不整合ポテンシャルによる調和振動子モデル――}
\author{}
\date{}

\begin{document}

\maketitle

\section{はじめに}

二重時空理論では、すべての質量を持つ粒子が内在的に通常時空と双対時空のペアを持つ。これらは双四元数代数 $\mathbb{H}\otimes\mathbb{H}\cong\mathrm{Cl}(3,1)$ により記述され、完全ローターは
\begin{equation}
R_{\text{total}} = R_{\text{usual}} R_{\text{dual}}
= \exp \left(\frac{\theta}{2} \hat{p} + \frac{\phi}{2}\hat{q}\right)
\end{equation}
と表される。ここで $\hat{p} = p_1 iI + p_2 iJ + p_3 iK, ~ \hat{q} = q_1 I + q_2 J + q_3 K$ であり、$\hat{p}$ はブーストベルソル（平方 $+1$）、$\hat{q}$ は双対セクターの回転ベルソル（平方 $-1$）である。

重力は相対ローター $\Omega = R_{\text{usual}}^\dagger R_{\text{dual}}$ の対数から得られるバイベクトル $\Omega_{\text{biv}}=\log\Omega$ のキリング形式
\begin{equation}
J = \frac{1}{16} B(\Omega_{\text{biv}},\Omega_{\text{biv}})
\end{equation}
により記述され、これはテレパラレル重力（TEGR）と古典的に等価である。

本節では、このねじれ不整合エネルギー $J$ を個別粒子レベルで内在的ポテンシャルとして採用し、変分原理から双対ローター遅れの厳密な動態方程式を導出する。これにより、質量による遅れが自然に現れ、自由粒子のド・ブロイ位相が再現されることを示す。

\section{粒子内在行動原理}

個別粒子の内在的動態を記述する有効行動を次のように提案する：
\begin{equation}
S_{\text{particle}} = \int \left[ 
\frac{1}{2} \dot{\phi}^2 \, \hat{p}\cdot\hat{p}
- \frac{1}{2} \dot{\theta}^2 \, \hat{q}\cdot\hat{q}
- \frac{m}{2} (\theta - \phi)^2 
\right] dt
\end{equation}
ここで、
\begin{itemize}
\item $\frac{1}{2}\dot{\theta}^2$：通常セクターの運動エネルギー（ブースト自由度）
\item $-\frac{1}{2}\dot{\phi}^2$：双対セクターの符号反転（時間反転性を反映）
\item $\frac{m}{2}(\theta - \phi)^2$：ねじれ不整合ポテンシャル（$J$ に比例）
\item $m$：粒子の質量、ポテンシャルの剛性を表す
\end{itemize}
である。この行動は、多数の粒子の集団平均で古典重力の $\int J\,d^4x$ を再現する。

\section{運動方程式の導出}

ラグランジアンを
\begin{equation}
L = \frac{1}{2} \dot{\phi}^2 \, \hat{p}\cdot\hat{p}
- \frac{1}{2} \dot{\theta}^2 \, \hat{q}\cdot\hat{q}
- \frac{m}{2} (\theta - \phi)^2
\end{equation}
とする。オイラー・ラグランジュ方程式を適用する。

\begin{align}
\frac{d}{dt}\frac{\partial L}{\partial \dot{\theta}_a} &= \frac{\partial L}{\partial \theta_a}
\;\Rightarrow\; \ddot{\theta}_a - m (\theta_a - \phi_a) = 0,\\[6pt]
\frac{d}{dt}\frac{\partial L}{\partial \dot{\phi}_a} &= \frac{\partial L}{\partial \phi_a}
\;\Rightarrow\; -\ddot{\phi}_a - m (\theta_a - \phi_a) = 0.
\end{align}
ここで、$\theta_a$ と $\phi_a$ はそれぞれ通常および双対セクターの独立した空間方向の成分である。

ねじれ不整合角 $\delta\theta_a = \theta_a - \phi_a$ を導入すると、
\begin{equation}
\ddot{\delta\theta}_a + m \delta\theta_a = 0.
\end{equation}
これは角振動数 $\sqrt{m}$ の調和振動子方程式である。

\section{自由粒子の解とド・ブロイ位相}

自由粒子（外部加速なし）の場合、$\dot{\theta}_a = \gamma v_a/c$（定数）と仮定する。初期条件を静止系で同期（$\delta\theta_a(0)=0$）とし、初期遅れ率を考慮すると、一般解は
\begin{equation}
\delta\theta_a(t) = A \sin(\sqrt{m}\, t) + B \cos(\sqrt{m}\, t).
\end{equation}
コンプトンスケールでの高周波振動を平均化する低エネルギー近似（$\sqrt{m}\, t \ll 1$）では線形蓄積
\[
\delta\theta_a(t) \approx \dot{\delta\theta}_a(0)\, t
\]
が支配的となる。ここに $\dot{\delta\theta}_a(0) \propto \gamma v_a$（質量による遅れ初期率）が入る。

位相として積分
\[
\int \delta\theta_a\, dt \propto (\gamma m c^2 t - \gamma m v_a x)/\hbar
\]
（固有時から座標時への変換後）を得る。これは正にド・ブロイ位相
\[
\frac{E t - \mathbf{p}\cdot\mathbf{x}}{\hbar}
\]
に一致する。

\section{量子近似への接続}

コンプトンスケールでの急速振動を平均化すると、$\delta\theta_a$ のゆっくり変化成分のみが残り、$\phi_a \approx -\theta_a +$（遅れ項）となる。この遅れ項の線形進化が波動関数の位相 $S(x,t)/\hbar$ に帰着し、低エネルギー有効理論としてシュレーディンガー方程式が自然に現れる。

\section{結論}

以上により、古典重力のキリング形式 $J$ を粒子内在ポテンシャルとして用いることで、変分原理から双対ローター遅れの厳密な動態方程式
\[
\ddot{\delta\theta}_a + m \delta\theta_a = 0
\]
を導出できた。このモデルは質量 $m$ をポテンシャル剛性として自然に取り入れ、自由粒子におけるド・ブロイ位相を正確に再現する。二重時空理論の双対ローター遅れ機構として最も整合的かつ背景独立な枠組みを提供するものである。

\end{document}